\scanpage{303}
Tribu 43

NOTES SUR LES SUPPLEMENTS DE THEORIE DES CORPS

(Rédaction n$^\circ$ 267, Tribu 43, Marlotte, Octobre 57)

I. - \uline{Matières devant aller en Alg. IV-V}, \uline{2ème édition}.

Le lemme 3, p.12 (Alg.V).

\uline{Suppléments sur les dérivations} (Extraits du \S3).

Le n$^\circ$1 de rappels est inutile.

Le début du n$^\circ$2 (jusqu'à la prop.1 exclue), où on fabrique les différentielles de Kähler, ira au chap.IV. Regarder les rédactions précédentes (projet Koszul de réédition d'Alg.II-III, diplodocus Grothendieck) où le sujet est traité. On dira ``\unsure{diffelles|diffélles} de Kähler'', car on ne sait pas bien quelle est la bonne notion de ``différentielle tout court'' (Dans bien des cas on doit chatrer celles de Kähler). La remarque p.18 mérite mieux qu'une remarque. Vers la l.10 de la p.19, il faut supposer que la base de M est finie. Donner aussi la suite exacte analogue de (1) pour les modules de différentielles.

A partir de la prop.1, ça va au chap.V. Dans la prop.1 on se ramène au cas d'ext. monogènes de 3 types particuliers (cf.n$^\circ$1). Le lemme 1 s'appellera ``lemme d'Audin''. Il y a des fautes de frappe à la p.20 (B et A ont été plusieurs fois intervertis).

Au n$^\circ$ 3, on commence par ne pas parler de dérivations. Notion de famille p-indépendante (linéaire indép. des monômes qu'il faut) ; ça veut dire qu'aucun élément n'est dans le sous corps engendré par les autres. Zornification ; existence de p-bases ; ce sont les familles p-indépendantes qui engendrent l'extension. Pour que la p-base ait d-éléments, il faut et il suffit que le degré de l'extension soit $p^d$ ;

\mspage{302}
